\section{Checking the scope of privacy claims}
\label{sec:scopechecks}

\subsection{Attack success under broader access}

Observer privileges determine which attack results transfer. The following
observation distinguishes an available successful attack from a tested
attack that finds no signal.

\begin{samepage}
\begin{proposition}[Transfer under nested observations]
Fix the target distribution, success measure and resource allowance.
Suppose observer $B$ receives everything available to observer $A$, and
$B$'s allowed attack class contains every admissible attack for $A$
composed with projection to $A$'s view, within the stated resource
allowance. Then $B$ can achieve at least the success of any particular
attack available to $A$. Non-detection by a tested attack for $A$ supplies
no upper bound on the success available to $B$.
\end{proposition}
\begin{proof}
Observer $B$ projects its observations to $A$'s view and executes the same
attack. Its output distribution and success remain unchanged. The
construction shows that an attack is available. A particular implementation
or optimiser may still fail to discover it, and failure by a tested attack
supplies no upper bound.
\end{proof}
\end{samepage}

The shared-key orthogonal-wrapper experiment illustrates this distinction.
It protects against an evaluated passive wire decoder while an authorised
cloud restores the original computation. For orthogonal $Q$,
\begin{equation}
 h'=hQ,\qquad h=h'Q^{\mathsf T}.
 \label{eq:rotation}
\end{equation}
Key possession changes the observer's view. A cloud that holds or derives
$Q$ falls outside the wire observer's no-key assumption, and a compromised
host can inspect that key and invert the transform while remaining passive.
The source account and later collection manifest expressly scope the
restoration result to the wire and logged captures (\artifact{E8},
\artifact{ER}). The measured result remains valid there, while
confidentiality from the restoring host requires a different argument.

Orthogonal wrapping also preserves norms and inner products. An invertible
coordinate change leaves information recovery unresolved when a fixed
decoder fails. The result describes the tested decoder. It provides
no universal division between structure and statistics, or upper bound on
information retained in the transcript. The robustness record itself
includes stronger labelled-pair attackers and depth-dependent recovery
(\artifact{E8}).

\begin{table}[t]
\centering\small
\begin{tabular}{@{}rrrr@{}}
\toprule
Depth & Estimated $K_{50}$ & Measured bracket & $K_{50}/5120$ \\
\midrule
1 & 258.32 & 200 to 300 & 0.05045 \\
4 & 355.81 & 300 to 400 & 0.06949 \\
8 & 617.93 & 600 to 700 & 0.12069 \\
\bottomrule
\end{tabular}
\caption{\textbf{Recovery crossings vary with depth at fixed width.}
The 27B/$H=5120$ artifact estimates absolute 50\% token-recovery crossings
by interpolation between measured $K$ values. Uncertainty in these
conditional crossing estimates was not quantified.
The 2.39-fold depth spread at fixed width does not establish a common
$0.09H$ coefficient. The source is \artifact{E8}.}
\label{tab:k50}
\end{table}

The pooled depth-1 session-rotation score is $3.71\%$, while individual
sessions include $5.73\%$, $3.65\%$ and $2.34\%$. The distinction changes
what the score means. A pooled value describes the aggregate and cannot
serve as a worst-session confidentiality bound. The empirical crossings in Table~\ref{tab:k50}
depend on the attack. Useful partial recovery can occur below the absolute 50\%
crossing, so these laboratory curves alone cannot determine a safe key
lifetime (\artifact{E8}).

\subsection{Comparing rotation cadences}

The rotation deployment study separates a three-session Mistral-7B
comparison from a later eight-run Mistral-NeMo-12B per-epoch matrix.
Fixed and per-session rotated keys have negligible measured throughput
differences in the first comparison. In the later matrix, identity is
correct in 20 of 24 sessions and four of eight runs are admitted with known
limitations. The N=256 setting advances no epochs. Its short request horizon makes
it a no-rotation control. Wire attacks on rotated epochs
record zero top-1 recovery within the tested decoder and capture scope
(\artifact{ER}).

The faster rotation carries a measured cost. For N=16, RTT=0 and session~1,
the identity record contains $2.8$ tokens/s for the ratchet and $7.6$
tokens/s for the static arm, a $63.16\%$ throughput reduction at the
recorded precision. The ledger reports about $5.5$ to $5.6$ tokens/s for
the evidence-triggered budget arm. We retain run and session
timings separately. Their extraction adds no new paired runtime estimate,
and models, cadences and RTT twins remain separate observations rather
than independent pooled replications (\artifact{ER}).

The historical records require the same care. An earlier narrative marks
the deployed matrix invalid, whereas a later manifest admits supporting
results after further analysis. The later disposition governs our reading.
Its identity, runtime and observer limitations remain attached to those
results, preserving both the revised assessment and the history from
which it arose.

\subsection{Reading mutual information estimates}

The MI-estimator analysis records near-zero estimates on some defended
boundaries (\artifact{MI}). At the population level, the variational
objective of Donsker and Varadhan for a critic $T$ satisfies
\begin{equation}
 L_T=\mathbb E_{P_{XZ}}[T(X,Z)]-
 \log\mathbb E_{P_XP_Z}[e^{T(X,Z)}]\le \operatorname{MI}(X,Z).
 \label{eq:dv}
\end{equation}
Here $\operatorname{MI}(X,Z)$ denotes mutual information. The bound points
in one direction. Finite-sample estimation can be difficult
\citep{song2020,mcallester2020}, and the population inequality itself
cannot supply an upper privacy budget.

Let $X$ be uniform on $m$ symbols and let $Z=X$. A constant critic $T=0$
has $L_T=0$, while $\operatorname{MI}(X,Z)=\log m$ nats. With $m=256$,
the channel reveals eight bits despite the zero critic readout. The example
is purely analytic. It demonstrates the inequality's direction, without
constituting a synthetic benchmark or estimating the defended system's
actual leakage. Even a valid zero lower bound leaves a small upper bound
unestablished. A second weak lower-bound critic supplies the same limited
inference, while an empirical plug-in estimate also needs statistical
justification before it can be called a certified population bound. The
source's own limitations state that neither estimator is an upper bound
(\artifact{MI}).

We therefore retain the estimator-specific observations while withholding
a small-leakage certificate. The inequalities alone leave actual mutual
information unresolved. They also cannot establish high leakage in every
defended cell. Capacity claims need comparable care.
Similar results across finitely tested architectures and schedules show
empirical saturation. An information-theoretic ceiling requires a bound
over all remote computations.
